\documentclass[12pt]{article}
\usepackage[utf8]{inputenc}
\usepackage[T1]{fontenc}
\usepackage{hyperref}
\title{Central Tharu Language Resources for K-12 Curriculum Development}
\author{Compiled by Claude Sonnet 5 (claude-sonnet-5) for Andr\'as Kornai}
\date{}
\begin{document}
\maketitle

\section{Introduction}
Central Tharu (ISO 639-3: the; Glottocode chit1274), also called Madhya-Purbiya Tharu, is an Indo-Aryan language of Nepal and India with about 290,000 speakers. A web search found very little material specific to Central Tharu. Most results concern other Tharu varieties (Rana Tharu, Dangaura Tharu), which are distinct and should not be treated as Central Tharu resources.

\section{Primary Educational Resources}
No Central Tharu primer or textbook was found in search results. A Cambridge chapter describes bilingual Dangaura Tharu--Nepali textbooks used in grades I--III at one school. That is a different variety, so it is useful only as a model:
\url{https://www.cambridge.org/core/books/abs/simultaneous-identities/language-education-and-knowledgemaking/489537F62DFD4A9BB5FEF7592C8C96C0}

\section{Linguistic and Reference Materials}
A Rapid Word Collection workshop, co-sponsored by Mother Tongue Centre Nepal and SIL Nepal, was held in Kathmandu in September 2014 as the first step toward a Madhya-Purbiya Tharu dictionary. The search did not show whether a finished dictionary followed:
\url{https://www.webonary.work/2014/11/13/madhya-purbiya-tharu-speakers-begin-the-first-dictionary-for-their-language/}

\section{Religious and Community Resources}
None found in the search results.

\section{Digital and Audio Resources}
None found in the search results.

\section{Curriculum Development Context}
A World Vision page notes that mother-tongue schooling is a constitutional right in Nepal (2015):
\url{https://www.wvi.org/bs/node/122531}
Tharu is reported to be taught as a mother tongue at primary level in government-aided schools in Kailali district, but the variety taught is not confirmed. The available evidence for Central Tharu is thin, and further searching under ``Madhya-Purbiya Tharu'' is recommended.
\end{document}